\section{Introduction}
\label{sec:intro}

Teams validate AI judges at scale by their correlation with human ratings~\citep{liu2023geval,zheng2023judging}.
Pearson correlation measures the linear association between the judge's scores and human ratings, and rank correlations measure agreement in order.
Neither says whether the judge predicts how experts rate each input.
That distinction matters whenever a decision uses probabilities, thresholds or averages of the judge's output.

We model the judge as a conditional distribution.
Given an item $x$ (a problem, a candidate answer, a reference), expert ratings $r\in\{0,\dots,K-1\}$ vary across raters and occasions.
We write $p(r\mid x)$ for that distribution and $q(r\mid x)$ for the judge's forecast.
Fitting $q$ is ordinary supervised learning with a proper loss, and evaluating it is forecast verification~\citep{gneiting2007strictly}.
This memo supports the blog post that makes this argument~\citep{miranda2026bitter}, in the spirit of the bitter lesson~\citep{sutton2019bitter}: use the general method, give it data and measure how error falls with data.

Correlation cannot identify $p$.
Pearson correlation across items depends on a forecast only through its conditional means, and only up to a positive affine map.
On $K$ categories, $K-2$ shape directions per item stay free, and so do the level and the scale of the means (Section~\ref{sec:identification}, \appref{app:proofs}, \appref{app:moments}).
Identification does not settle the matter, because a flexible model trained on correlation can still land on the right distribution.
We therefore ask an empirical question.
How many human ratings does a cross-entropy-fitted judge need before it beats a correlation-fitted judge on held-out ratings, and does the answer survive across datasets, models and objectives?

This memo consolidates every experiment in the repository (30 active folders and four archived ones, numbered 00--33; \appref{app:index}) and builds on the identification preprint~\citep{miranda2026correlation}.
We make five contributions.
\begin{enumerate}
\item \textbf{Identification.} We prove what Pearson correlation, mean squared error, cross-entropy and the ranked probability score identify, bound moment recovery by cross-entropy, and give an exact three-item example with equal means and different distributions (Section~\ref{sec:identification}; derivations in \appref{app:moments}).
\item \textbf{Controlled studies.} On synthetic ratings with known distributions, cross-entropy training beats Pearson training from 12 ratings when the mean carries the signal (experiment 08, three ratings per task) and from 48 in a second toy whose spread and shape also vary (experiment 20, where it loses at 12).
When the mean says little, it needs 192 ratings against Pearson, Spearman and mixed objectives and 24 against Kendall (experiment 06, scenario A; Section~\ref{sec:controlled}).
\item \textbf{Frozen judges on human ratings.} With five frontier judges and a 27-setting calibration head, cross-entropy fitting clears Pearson fitting and a uniform forecast after 4--8 ratings, one per item, in all 25 dataset-by-judge cells of ChaosNLI and DICES (Section~\ref{sec:frozen}).
\item \textbf{Trained judges.} In matched LoRA fine-tuning with four objectives, Pearson-trained judges have higher raw RPS than CE-trained judges in all seven setups at the largest budget.
Squared error and CE differ by $0.0003$ RPS at $K=2$ (interval includes zero), and CE beats squared error on five of six panels with $K\ge3$ (Section~\ref{sec:lora}).
\item \textbf{Scaling laws.} A floor plus $aN^{-b}$ fits the held-out RPS curves of cross-entropy-fit judges in 9 of 11 panels (median $b=0.66$ where the floor is identified) and those of Pearson-fit judges in 1 of 11 (Section~\ref{sec:scaling}).
\end{enumerate}

\paragraph{Related work.}
\citet{han2025judgesverdict} report, across 54 judges, that high correlation with humans does not guarantee human-like agreement.
\citet{guerdan2025validating} show that validating judges against forced-choice human labels can select judges performing as much as 31\% worse than selection that accounts for rating indeterminacy.
\citet{elangovan2025beyond} show that one aggregate correlation can hide differences when human labels are uncertain.
We add training and a matched-budget comparison of objectives.
